\scanpage{60}
d'abord construire un $u:B_X\longrightarrow A_Y$ ayant cette dernière vertu, indépendamment de considérations cohomologiques. \deleted{Soit un diviseur sur $X\times B_X$}\inserted{Notons d'abord qu'il suffit de} faire le calcul avec les classes d'équivalence linéaire $z$, $d$ des cycles (ou $\unsure{\overline{\mathrm m}|=}$ l'équivalence d'Albanese) -- puisque si $\unsure{W|\Omega|\mathfrak{A}}\in\operatorname{Cyc}_0(Y)$ est lin. équi. à $0$, $S_Y(W)$ est nul. Donc on n'a pas besoin de se préoccuper de la question \deleted{que les produits}\inserted{de \deleted{\unsure{de}} l'existence} des \unsure{expressions}, et se borner à regarder $d\longmapsto S_Y(z(d))$. Or le calcul des intersections prouve que, si $\pi:Y\times X\times B_X\longrightarrow Y\times B_X$ est la projection canonique, \inserted{et $\Theta\in A^1(X\times B_X)$ l'élément canonique,} on a \deleted{(identifiant $d$ à un point}
\[
z(\Theta(b))\simeq\deleted{\unsure{\pi_*}}C(b)
\]
où $C\in Z^m(Y\times B_X)$ est donné par
\[
C=\pi_*\bigl((1_Y\times\Theta)\mathbin{.}(z\times1_{B_X})\bigr)
\]
\inserted{D'ailleurs $z(\Theta(0))=0$ donc $z(\Theta(b))=z(\Theta((b)-(0)))$}

On applique alors 1.1. à $C$ sur $Y\times B_X$, pour obtenir $u:A_{B_X}=B_X\longrightarrow A_Y$, satisfaisant
\[
\begin{aligned}
&\deleted{\unsure{u\,z(\Theta(b))|u\,\mathfrak{z}(\Theta(b))}}\\
&S_Y(C((b)-(0)))=u(b)\qquad\text{i.e.}\\
&S_Y(z(\Theta(b)))=u(b).
\end{aligned}
\]
% archive sheet number: 59

